\section{Etat de l'art}
Cette partie propose une analyse des méthodes de l'état de l'art pour expliquer les graphes. 
Elle commence par détailler le rôle de la prédiction de lien dans les graphes, ainsi les méthodes existantes.
La seconde propose un apperçu des travaux se concentrant sur l'inférence de features débiaisées ou décorrélées. 
La troisième détaille les pistes d'explicabilité utilisées dans l'état de l'art.


\subsection{Algorithmes de prédiction de lien}
Les graphes sont des structures mathématiques définies par leurs nœuds et leurs arêtes. 
Formellement, un graphe est représenté par un quadruplet $(V, E, \chi, \chi_E)$, où $V$ 
désigne l'ensemble des nœuds, $E \subseteq V \times V$ représente l'ensemble des arêtes, 
et $\chi : V \rightarrow \Omega_V$ et $\chi_E : (V \times V) \rightarrow \Omega_E$ correspondent respectivement aux attributs des nœuds et aux attributs des paires de nœuds.

La prédiction de lien consiste à utiliser ces informations $(V, E, \chi, \chi_E)$ afin de prédire de nouveaux éléments de $E$. Elle s'appuie notamment sur de nombreux algorithmes permettant d'enrichir $\chi$ ou $\chi_E$ avec de nouvelles informations. L'intérêt de la prédiction de lien est multiple : 
\begin{itemize}
    \item La plupart des réseaux réels ne sont observés que de manière incomplète. Les algorithmes capables de prédire avec précision quels liens sont manquants peuvent accélérer considérablement la collecte de données sur les réseaux et améliorer la validité des modèles les utilisant.
    \item Elle permet de modéliser et de comprendre les mécanismes d'évolution et de dynamique des réseaux temporels, en anticipant la trajectoire future des connexions (par exemple, la formation de nouvelles collaborations scientifiques).
    \item Elle permet de nombreuses tâches utilisées dans l'industrie telles que la recommandation d'amis ou de contenus sur les réseaux sociaux, la suggestion de produits sur les plateformes de commerce électronique, ou encore la détection de fraudes et d'activités suspectes dans les réseaux de transactions financières. Elle trouve une application cruciale en bio-informatique pour la découverte biomédicale, notamment pour prédire des interactions inconnues entre protéines (PPI), des associations gène-maladie, ou pour le repositionnement de médicaments (drug repositioning) en modélisant les interactions molécules-cibles.
    \item La performance en prédiction de lien peut être utilisée comme un indicateur d'à quel point des métriques expliquent correctement la structure d'un graphe. C'est pour cette utilisation spécifique que ce projet l'utilise.
\end{itemize}

Un constat fondamental de la littérature est qu'aucun algorithme de prédiction de lien n'est universellement supérieur aux autres sur l'ensemble des topologies de graphes. Les performances de chaque approche dépendent intrinsèquement de la nature du réseau étudié. De manière générale, ces algorithmes peuvent être décomposés en trois grandes catégories :

\textbf{Les métriques topologiques :} Ces approches calculent des descripteurs déterministes. On distingue parmi elles des métriques \textit{globales} du réseau ($V$, $E$, degré moyen, variance de la distribution des degrés, diamètre, assortativité, transitivité ou coefficient de clustering global), des caractéristiques \textit{nodales} calculées indépendamment pour chaque sommet (degré, proximité, interstitialité, centralité de vecteur propre, Katz centralité, PageRank, charge, coefficient de clustering local ou volume de triangles locaux), et enfin des métriques \textit{de paire de noeuds}. Ces dernières exploitent la structure locale ou les chemins entre $i$ et $j$ via des indices de voisinage (voisins communs, Jaccard, Adamic/Adar, \textit{Resource Allocation}, katz-index\cite{katz1953}), la longueur du plus court chemin, le PageRank personnalisé, l'attachement préférentiel (produit des degrés), ou des approximations de rang faible (LRA) basées sur la décomposition en valeurs singulières (SVD) de la matrice d'adjacence, évaluées par produit scalaire ou par agrégation locale du voisinage. Une taxonomie presque exhaustive de ces prédicteurs est notamment utilisée par Ghasemian et al. \cite{01stackingLinkPred}.
    
\textbf{L'inférence de communautés :} Ces méthodes exploitent l'hypothèse d'homophilie structurelle selon laquelle les nœuds partagent une plus forte propension à se connecter s'ils appartiennent à des sous-groupes denses. Il en existe un multitude, plutôt qu'une liste exhaustive voici le détail de quelques unes des plus pertinentes pour ce projet.

L'algorithme de Louvain \cite{louvain2008}, qui repose sur une approche gloutonne heuristique visant à maximiser la \textit{modularité} du réseau, métrique mesurant l'écart de densité de liens internes aux communautés par rapport à un modèle nul (\textit{Null Model}) de configuration aléatoire. Il procède de manière itérative par optimisation locale de l'apport de modularité puis par contraction du graphe en super-nœuds. 

À l'inverse, le modèle en blocs stochastiques (SBM) \cite{sbm1983} adopte une démarche probabiliste et générative : les nœuds sont assignés à des variables latentes de blocs, et la probabilité de connexion entre deux nœuds dépend exclusivement d'une matrice de transition inter-groupes, permettant de capturer des structures non seulement communautaires, mais aussi disassortatives (biparties) ou de type centre-périphérie. 

Des approches alternatives comme \textit{Infomap} \cite{infomap2008} formalisent ce problème sous l'angle de la théorie de l'information \cite{infotheory2023} en minimisant la longueur de description du flux d'une marche aléatoire via le codage de Huffman.

\vspace{4px}

\textbf{L'inférence d'embeddings (Représentations géométriques latentes) :} Ces techniques projettent la topologie non linéaire du graphe dans un espace vectoriel continu de faible dimension $\mathbb{R}^d$ où les propriétés de proximité sont préservées. 

\textit{DeepWalk} \cite{deepwalk2014} génère des trajectoires de sommets via des marches aléatoires, puis optimise les représentations via le modèle \textit{Skip-gram} pour apprendre des représentations qui capturent le voisinage des nœuds.

\textit{node2vec} \cite{node2vec2016} généralise ce formalisme en introduisant des marches de second ordre biaisées par deux hyperparamètres $p$ (paramètre de retour) et $q$ (paramètre d'exploration). Ce mécanisme permet d'orienter les marches aléatoires selon une stratégie de recherche en largeur (BFS) ou en profondeur (DFS). Il a été théorisé que ces stratégies permettent respectivement de mieux capturer l'homophilie locale (à l'instar des communautés), ou le rôle global de chaque nœud (centralité).

Pour exploiter ces descripteurs hétérogènes, de nombreuses approches existent. L'approche historique consiste à concaténer l'ensemble des métriques topologiques et des caractéristiques nodales pour entraîner un classificateur tabulaire standard, tel que les forêts aléatoires ou \textit{XGBoost} \cite{lichtenwalter2010link}. C'est celle qui sera majoritairement utilisée dans ce projet de part sa faible complexité algorithmique.

L'approche de référence de Ghasemian et al. \cite{01stackingLinkPred}  combine les prédictions probabilistes issues de différents modèles au sein d'un méta-classificateur (\textit{Stacking}), démontrant qu'une telle intégration non linéaire permet d'atteindre une borne de prédictibilité quasi-optimale. Dans cette même lignée, de récents travaux introduisent des mécanismes de mélange d'experts pour agréger dynamiquement les prédictions selon le contexte topologique local \cite{ma2024mixture}.

Par ailleurs, la prédiction de liens est fréquemment formulée comme un problème d'algèbre linéaire via la complétion de matrices et la factorisation matricielle non négative (NMF) \cite{menon2011link}, ou encore sous la forme de modèles graphiques probabilistes (\textit{Probabilistic Graphical Models}), à l'instar des approches bayésiennes et des champs aléatoires de Markov \cite{taskar2003link}.

Enfin, l'état de l'art contemporain s'est massivement déplacé vers les réseaux de neurones sur graphes (\textit{Graph Neural Networks} - GNN) \cite{kipf2016semi}. Ces architectures apprennent de bout en bout (\textit{end-to-end}) des fonctions d'agrégation et de passage de messages (\textit{message passing}) au sein des voisinages locaux, court-circuitant l'ingénierie manuelle de caractéristiques pour optimiser directement une fonction de perte basée sur la reconstruction géométrique de la matrice d'adjacence.


\subsection{Factorisation orthogonale, débiaisage et modèles nuls}
\label{sec:decorrelation}

Afin d'isoler l'effet spécifique d'une caractéristique structurelle ou d'un attribut de nœud, il est nécessaire de s'affranchir des corrélations sous-jacentes. Dans la littérature, ce problème d'indépendance statistique s'articule autour de trois axes méthodologiques complémentaires.

Le premier axe concerne le \textbf{débiaisage des représentations latentes}. Initialement développé pour garantir l'équité temporelle ou sociologique, il introduit la notion de (\textit{fairness}). Des approches comme \textit{Fairwalk} \cite{rahman2019fairwalk} reprennent la principe développé par Deepwalk en modifiant l'étape de transition des marches aléatoires pour équilibrer la probabilité de saut entre différents groupes. Bien qu'initialement développé pour débiaiser les marches aléatoires entre différents groupes sociologiques (ex : hommes/femmes), il peut être étendu à d'autres ensembles discrets de groupes, telles que les communautés. \textit{CrossWalk} \cite{khajehnejad2022crosswalk} reprend également le mécanisme de Deepwalk et applique des poids de minimisation des disparités aux frontières des communautés. Si ces méthodes s'avèrent empiriquement efficaces pour décorréler les \textit{embeddings} d'attributs sensibles (sexe, origine), leurs formulations mathématiques reposent encore 
sur des choix largement heuristiques. À ma connaissance, aucune étude formelle ni comparative ne garantit à ce jour l'optimalité de ces fonctions d'agrégation, laissant ouverte la possibilité de formulations plus performantes.

Le deuxième axe formalise mathématiquement cette indépendance à travers l'\textbf{apprentissage de représentations démêlées} (\textit{disentangled representation learning}) et l'orthogonalité géométrique. Inspiré du traitement automatique des langues où l'élimination des biais repose sur la projection géométrique dans le complément orthogonal d'un sous-espace conceptuel \cite{schumann2024prejudice}, ce paradigme a été transposé aux graphes. Le modèle \textit{ADGCN} \cite{zheng2021adversarial} utilise ainsi des réseaux antagonistes (\textit{adversarial learning}) couplés à des pertes d'indépendance de distance afin de forcer la séparation de $K$  caractéristiques de nœuds dans des espaces latents distincts et perpendiculaires. Ces mécanismes permettent ainsi d'obtenir $K$ facteurs d'existance de lien non juxtaposéss, facilitant grandement l'explicabilité. A noter que l'interprétation de chacun de ces facteurs reste à la charge des utilisateurs.

Enfin, le troisième axe — dont s'inspire directement ce projet — s'appuie sur la théorie des \textbf{modèles nuls} pour neutraliser l'impact d'une structure dominante lors de l'évaluation d'une autre propriété. Plutôt que de projeter géométriquement des vecteurs, cette approche ajuste les fonctions d'évaluation statistique. Par exemple, Cazabet et al. \cite{cazabet2022enhancing} redéfinissent la fonction de modularité de Louvain en y injectant un modèle nul spatial contraint par les degrés, permettant de mettre en exergue de nouvelles communautés intéressantes que l'algorithme classique ne mettait pas en valeur. Cette interaction entre structures géométriques et probabilistes est d'ailleurs formalisée par Duvivier et al. \cite{duvivier2024graph} pour évaluer rigoureusement la pertinence statistique des modèles de graphes. C'est précisément cette logique de soustraction d'effets probabilistes de référence qui est exploitée dans ce projet pour formuler nos méthodes d'inférence de métriques décorrélées.


\subsection{Intelligence artificielle explicable (XAI) sur graphes}

L'avènement de modèles complexes et de boîtes noires pour la prédiction de liens a propulsé le développement de méthodes d'explicabilité (XAI). Dans la littérature, les méthodes de XAI appliquées aux structures de graphes se divisent principalement selon deux axes taxonomiques : la nature de l'explication (approches intrinsèques, où le modèle est nativement interprétable, versus approches \textit{post-hoc}, qui analysent un modèle déjà entraîné) et le niveau d'abstraction (explications locales centrées sur une arête ou un nœud, versus explications globales synthétisant le comportement du réseau entier).

La majorité des travaux récents dédiés aux graphes, à l'instar des approches explorées dans \cite{kamal2025game}, se focalisent sur la découverte de règles structurelles (\textit{rule discovery}) ou l'extraction de motifs sous-graphiques pertinents pour expliquer les prédictions de réseaux de neurones sur graphes (GNN). Cependant, ces formalismes sont spécifiques aux GNNs et s'appliquent difficilement à l'explicabilité d'autres modèles. Pour pallier cela, l'adaptation des méthodes d'attribution fondées sur la théorie des jeux coopératifs, et plus particulièrement les valeurs de Shapley \cite{lordShapleyOriginal_1953}, s'impose comme une alternative robuste grâce à ses garanties axiomatiques uniques (additivité, efficacité, symétrie et contribution nulle du joueur nul).

Le principal verrou scientifique de l'utilisation des valeurs de Shapley réside dans leur complexité computationnelle exponentielle, le calcul exact des contributions nécessitant l'évaluation de $2^M$ coalitions possibles pour un ensemble de $M$ variables. Pour lever cette limite, des alternatives théoriques récentes ont été proposées, à l'image des valeurs FESP (\textit{Fair and Efficient Shapley-based Permutations}) \cite{fesp2022}, qui optimisent le calcul d'attribution en réduisant l'espace des permutations admissibles tout en préservant certaines propriétés de justice. Néanmoins, dans le cadre de ce projet, les valeurs SHAP classiques ont été retenues au profit d'approches par échantillonnage de Monte-Carlo usuellement utilisées dans la littérature. Ce choix se justifie par la nécessité de préserver l'axiome d'additivité stricte à l'échelle locale, condition strictement nécessaire pour permettre une explicabilité globale par simple sommation des effets locaux.